\section{Introduction}
\label{sec:introduction}

Zero-shot anomaly detection aims to identify abnormal images or regions without
training examples from the target category. Vision-language models make this
setting practical by comparing visual features with language descriptions of
normal and abnormal states. WinCLIP~\cite{jeong2023winclip} introduced
compositional normal/anomalous prompts and local window scoring for
zero-/few-shot anomaly classification and segmentation. AnomalyCLIP~
\cite{zhou2024anomalyclip} further learns object-agnostic normality and
abnormality prompts from auxiliary anomaly data, and later methods improve the
same family of models with hybrid prompts~\cite{cao2024adaclip}, visual
context prompting~\cite{qu2024vcpclip}, and anomaly-aware adapters~
\cite{ma2025aaclip}. These works show that CLIP can provide strong
normal/abnormal semantic prototypes.

However, anomaly localization is not only a semantic matching problem. A patch
with a strong abnormal response is still a context-dependent observation. Local
texture changes, high-frequency responses, and boundary discontinuities often
appear near defects, but they also occur in normal materials, object contours,
and imaging noise. The same response may indicate a scratch on a smooth metal
surface, while being normal on a woven texture. Treating local response strength
as final evidence therefore mixes true defects with valid local variation.

This distinction is important because frequency and local-structure cues are
already used in CLIP-based anomaly detection. FE-CLIP injects frequency
information through trained adapters~\cite{gong2025feclip}; WMoE-CLIP uses
Haar wavelet features in prompt learning~\cite{chen2026wmoeclip}; HarmoniAD
models local structure and global semantics through frequency-domain branches~
\cite{zhang2026harmoniad}. Our claim is not that frequency is a new useful
signal. Instead, we ask where such local signals should enter the scoring
mechanism. The controlled results in this project indicate that direct
wavelet-map fusion is weak, while using boundary-aware local reliability to
select patch evidence and calibrate prototypes is more stable.

We therefore formulate CLIP-based zero-shot anomaly localization as
image-conditioned normal reference estimation. The method keeps the
CLIP/AnomalyCLIP representations fixed, uses the test image itself to select
reliable normal and abnormal patch evidence, and recomputes local anomaly
scores with image-conditioned prototypes. This reframes the score from "how
strong is the local response" to "how abnormal is this patch relative to the
normal reference estimated from the current image."

Our contributions are:
\begin{itemize}
  \item We identify context-dependent local evidence as a central difficulty in
  CLIP-based zero-shot anomaly localization: local responses are candidate
  cues, not anomaly evidence by themselves.
  \item We introduce an image-conditioned prototype calibration framework that
  estimates a per-image normal reference from reliable patch evidence without
  target-category samples, labels, or test-time parameter updates.
  \item We organize the current experiments as mechanism tests. Direct
  wavelet-map fusion serves as a negative control, semantic-only prototype
  adaptation is a strong baseline, and boundary-aware reliability plus
  conservative calibration gives the best controlled row on MVTec AD and VisA.
\end{itemize}

We deliberately avoid presenting the method as a ranking story. The
AnomalyCLIP backbone and prompts are learned from auxiliary anomaly data; the
claim here is narrower. The added module is a test-time evidence-selection and
prototype-calibration step on top of frozen representations.
